%! Diagonalizing a Matrix — Unit 6, Lesson 6.2 — Exit Ticket (Answer Key)
\documentclass[10pt]{article}
\usepackage{linalg-article}
\usepackage{linalg-key}   % pulls in -boxes; do NOT also load -boxes
\newcommand{\TallMath}[1]{$\displaystyle #1\rule[-1.4em]{0pt}{3.2em}$}
\newcommand{\xx}{\mathbf{x}}
\newcommand{\zero}{\mathbf{0}}

\begin{document}

\pageheader{Unit 6, Lesson 6.2}{Exit Ticket --- Answer Key}
\namedateperiod

\vspace{0.15in}

No notes. Show your reasoning. Use $A=\begin{bmatrix} 1 & 4 \\ 0 & 3 \end{bmatrix}$ for items 1--2.

\begin{enumerate}[leftmargin=1.8em, itemsep=14pt]
  \item \textbf{Build $X$ and $\Lambda$.} State the two eigenvalues (it is triangular), find an eigenvector
        for each, and write the eigenvector matrix $X$ and the eigenvalue matrix $\Lambda$.
        \ansline{Triangular $\Rightarrow\lambda=1,3$. $\lambda=1$: $A-I=\begin{bsmallmatrix}0&4\\0&2\end{bsmallmatrix}\Rightarrow x_2=0\Rightarrow\begin{bsmallmatrix}1\\0\end{bsmallmatrix}$; $\lambda=3$: $A-3I=\begin{bsmallmatrix}-2&4\\0&0\end{bsmallmatrix}\Rightarrow x_1=2x_2\Rightarrow\begin{bsmallmatrix}2\\1\end{bsmallmatrix}$. So $X=\begin{bsmallmatrix}1&2\\0&1\end{bsmallmatrix}$, $\Lambda=\begin{bsmallmatrix}1&0\\0&3\end{bsmallmatrix}$.}
  \item \textbf{Diagonalize.} Find $X^{-1}$ (Unit~2 formula) and write $A$ as $X\Lambda X^{-1}$.
        \ansline{$\det X=1$, so $X^{-1}=\begin{bsmallmatrix}1&-2\\0&1\end{bsmallmatrix}$. Then $A=X\Lambda X^{-1}=\begin{bsmallmatrix}1&2\\0&1\end{bsmallmatrix}\begin{bsmallmatrix}1&0\\0&3\end{bsmallmatrix}\begin{bsmallmatrix}1&-2\\0&1\end{bsmallmatrix}=\begin{bsmallmatrix}1&6\\0&3\end{bsmallmatrix}\begin{bsmallmatrix}1&-2\\0&1\end{bsmallmatrix}=\begin{bsmallmatrix}1&4\\0&3\end{bsmallmatrix}.\ \checkmark$}
  \item \textbf{Explain.} Once you have $A=X\Lambda X^{-1}$, why is $A^{k}=X\Lambda^{k}X^{-1}$ so much easier
        to compute than multiplying $A$ by itself $k$ times? (One or two sentences.)
        \ansline{Every inner $X^{-1}X$ collapses to $I$, leaving $A^k=X\Lambda^k X^{-1}$; and $\Lambda$ is diagonal, so $\Lambda^k$ is just each eigenvalue raised to the $k$ --- no repeated matrix multiplication.}
\end{enumerate}

\begin{teachernote}
Sort into ``got it / arithmetic slip / concept slip.'' Item~1 checks lining eigenvectors into $X$ and
eigenvalues into $\Lambda$ ($\lambda=1,3$; $X=\begin{bsmallmatrix}1&2\\0&1\end{bsmallmatrix}$); item~2 is the
diagonalization with a \emph{fraction-free} inverse ($\det X=1$); item~3 is the target idea --- the
$X^{-1}X$ cancellation plus $\Lambda$ diagonal. Watch for column/eigenvalue order mismatches and for
students who invert $X$ incorrectly. If many miss item~3, re-anchor on the $A^2=X\Lambda^2X^{-1}$
derivation from the notes before Lesson~6.3.
\end{teachernote}

\end{document}
